% ============================================================
% UFUU Bell Paper — T4 Uniqueness Supplement
% File: UFUU_Bell_T4_Uniqueness_20260502.tex
% Author: W. Jason Tuttle (ORCID: 0009-0003-3830-0551)
% Date: 2026-05-02
% Status: Verified computationally (Python/NumPy/SciPy, 2026-05-02)
%
% Purpose: Standalone LaTeX for Section 3.3 of the Bell paper.
%   Proves that Tuttle's Triad T1-T3, augmented by the minimum
%   circuit depth constraint T4, uniquely selects the CNOT
%   Weyl class (up to local unitaries) as the two-qubit
%   realization of the Recursive Binary Fold F.
%
% To include in the Bell paper:
%   % ============================================================
% UFUU Bell Paper — T4 Uniqueness Supplement
% File: UFUU_Bell_T4_Uniqueness_20260502.tex
% Author: W. Jason Tuttle (ORCID: 0009-0003-3830-0551)
% Date: 2026-05-02
% Status: Verified computationally (Python/NumPy/SciPy, 2026-05-02)
%
% Purpose: Standalone LaTeX for Section 3.3 of the Bell paper.
%   Proves that Tuttle's Triad T1-T3, augmented by the minimum
%   circuit depth constraint T4, uniquely selects the CNOT
%   Weyl class (up to local unitaries) as the two-qubit
%   realization of the Recursive Binary Fold F.
%
% To include in the Bell paper:
%   % ============================================================
% UFUU Bell Paper — T4 Uniqueness Supplement
% File: UFUU_Bell_T4_Uniqueness_20260502.tex
% Author: W. Jason Tuttle (ORCID: 0009-0003-3830-0551)
% Date: 2026-05-02
% Status: Verified computationally (Python/NumPy/SciPy, 2026-05-02)
%
% Purpose: Standalone LaTeX for Section 3.3 of the Bell paper.
%   Proves that Tuttle's Triad T1-T3, augmented by the minimum
%   circuit depth constraint T4, uniquely selects the CNOT
%   Weyl class (up to local unitaries) as the two-qubit
%   realization of the Recursive Binary Fold F.
%
% To include in the Bell paper:
%   % ============================================================
% UFUU Bell Paper — T4 Uniqueness Supplement
% File: UFUU_Bell_T4_Uniqueness_20260502.tex
% Author: W. Jason Tuttle (ORCID: 0009-0003-3830-0551)
% Date: 2026-05-02
% Status: Verified computationally (Python/NumPy/SciPy, 2026-05-02)
%
% Purpose: Standalone LaTeX for Section 3.3 of the Bell paper.
%   Proves that Tuttle's Triad T1-T3, augmented by the minimum
%   circuit depth constraint T4, uniquely selects the CNOT
%   Weyl class (up to local unitaries) as the two-qubit
%   realization of the Recursive Binary Fold F.
%
% To include in the Bell paper:
%   \input{UFUU_Bell_T4_Uniqueness_20260502}
%   or copy Section 3.3 verbatim into the main .tex file.
%
% Correction note (2026-05-02):
%   Earlier draft incorrectly claimed ZZ at pi/4 has zero
%   concurrence. ZZ IS a perfect entangler but requires circuit
%   depth 3 from |00> (two single-qubit prep gates + ZZ).
%   T4 is therefore correctly stated as minimum CIRCUIT DEPTH
%   from product input |00>, not minimum interaction rank alone.
%   The uniqueness result is unchanged; only the proof path differs.
% ============================================================

% ── If compiling standalone, uncomment the block below ──────
% \documentclass[12pt]{article}
% \usepackage{amsmath,amssymb,amsthm}
% \usepackage{physics}
% \usepackage{braket}
% \usepackage{hyperref}
% \newtheorem{theorem}{Theorem}
% \newtheorem{lemma}{Lemma}
% \newtheorem{definition}{Definition}
% \newtheorem{corollary}{Corollary}
% \begin{document}
% ── End standalone preamble ─────────────────────────────────

\subsection{Uniqueness of the Two-Qubit Realization (T4)}
\label{sec:T4_uniqueness}

Section~\ref{sec:fold_circuit} identified
$F_{\mathrm{circuit}} = \mathrm{CNOT} \cdot (H \otimes I)$
as the minimal unitary realization of the abstract fold operation $F$
in the two-qubit Hilbert space $\mathcal{H} = \mathbb{C}^2 \otimes \mathbb{C}^2$.
We now prove that this realization is \emph{unique} by introducing a
fourth structural constraint (T4) and showing that T1--T4 jointly select
the CNOT Weyl class as the only admissible realization, up to local unitaries.

% ── Weyl Chamber Background ──────────────────────────────────
\subsubsection*{Background: The Weyl Chamber}

Two-qubit gates in $SU(4)$ are classified up to local (single-qubit)
unitaries by their Weyl decomposition~\cite{Zhang2003}:
\begin{equation}
    U = (A_1 \otimes A_2)\,
        \exp\!\bigl[i(c_1\,XX + c_2\,YY + c_3\,ZZ)\bigr]\,
        (A_3 \otimes A_4),
    \label{eq:weyl_decomp}
\end{equation}
where $A_k \in SU(2)$, $XX = \sigma_x\otimes\sigma_x$,
$YY = \sigma_y\otimes\sigma_y$, $ZZ = \sigma_z\otimes\sigma_z$,
and $(c_1,c_2,c_3) \in [0,\pi/4]^3$ with $c_1 \geq c_2 \geq c_3 \geq 0$
are the Weyl coordinates. The interaction rank $\mathrm{rank}(U) =
\#\{k: c_k \neq 0\}$ counts the number of active interaction axes.
Key classes and their coordinates are given in Table~\ref{tab:weyl}.

\begin{table}[h]
\centering
\caption{Weyl coordinates and interaction rank for representative two-qubit gates.}
\label{tab:weyl}
\begin{tabular}{lcccc}
\hline
Gate class       & $c_1$    & $c_2$    & $c_3$    & Rank \\
\hline
Identity         & $0$      & $0$      & $0$      & 0 \\
CNOT             & $\pi/4$  & $0$      & $0$      & 1 \\
$\sqrt{i\mathrm{SWAP}}$ & $\pi/8$ & $\pi/8$ & $0$ & 2 \\
$i\mathrm{SWAP}$ & $\pi/4$  & $\pi/4$  & $0$      & 2 \\
B gate           & $\pi/4$  & $\pi/8$  & $0$      & 2 \\
$\sqrt{\mathrm{SWAP}}$ & $\pi/8$ & $\pi/8$ & $\pi/8$ & 3 \\
SWAP             & $\pi/4$  & $\pi/4$  & $\pi/4$  & 3 \\
\hline
\end{tabular}
\end{table}

% ── T4 Definition ────────────────────────────────────────────
\subsubsection*{The Fourth Constraint: Minimum Circuit Depth from Product Input (T4)}

\begin{definition}[T4 --- Minimum Circuit Depth]
\label{def:T4}
\textbf{(T4) Minimum circuit depth:}
The fold operation $F$ selects, from among all unitary realizations
satisfying T1--T3 and capable of producing maximal entanglement from
the unresolved binary product state $\ket{00}$, the realization requiring
the fewest elementary gates (single-qubit unitaries plus two-qubit
entangling gates) to achieve concurrence $\mathcal{C} = 1$ starting
from $\ket{00}$.
\end{definition}

The physical motivation is immediate. The recursive binary fold
$U = F(U,U)$ begins each step with two independent, unresolved binary
distinctions — represented as the product state $\ket{00}$. The
information-maximization primitive requires that maximal entanglement
be generated with minimum additional structure: the simplest circuit
that accomplishes this is the one the fold selects. T4 makes this
minimality requirement precise and measurable.

% ── Main Theorem ─────────────────────────────────────────────
\subsubsection*{Uniqueness Theorem}

\begin{theorem}[T1--T4 Uniquely Select the CNOT Class]
\label{thm:CNOT_uniqueness}
Among all two-qubit unitary gates, T1 (non-commutativity),
T2 (entanglement generation), T3 (partial irreversibility at
the recursive level), and T4 (minimum circuit depth from $\ket{00}$)
jointly select the CNOT Weyl class --- gates with Weyl coordinates
$(\pi/4, 0, 0)$ --- as the unique admissible realization of $F$,
up to local unitaries.
\end{theorem}

\begin{proof}
We establish the minimum achievable circuit depth (T4), then show
only the CNOT class achieves it while satisfying T1--T3.

\medskip
\noindent\textbf{Step 1: Lower bound on circuit depth.}
Any circuit producing maximal entanglement from $\ket{00}$ must
include at least one two-qubit entangling gate (since product
unitaries preserve the product structure of $\ket{00}$, violating T2)
and at least one single-qubit gate to create the superposition that
the entangling gate then correlates. The minimum circuit depth from
$\ket{00}$ to a maximally entangled state is therefore $d^* = 2$.

\medskip
\noindent\textbf{Step 2: Which rank-1 gates achieve $d^* = 2$ from $\ket{00}$?}
A depth-2 circuit from $\ket{00}$ has the form
$(A \otimes I) \cdot E$ or $(I \otimes A) \cdot E$ where $A \in SU(2)$
is a single-qubit gate and $E$ is a two-qubit entangling gate.
We require $\mathcal{C}(E\cdot(A\otimes I)\ket{00}) = 1$.

For the three single-axis interaction classes at $c = \pi/4$:

\begin{itemize}

\item \textbf{$XX$ axis}: $\exp(i\,\frac{\pi}{4}\,XX)$ with $A = H$
  (Hadamard) gives $\mathrm{CNOT}\cdot(H\otimes I)\ket{00} = \ket{\Phi^+}$,
  concurrence $= 1$. Depth 2 achieved. (Numerically verified:
  $\mathcal{C} = 1.000000$.)

\item \textbf{$YY$ axis}: $\exp(i\,\frac{\pi}{4}\,YY)$ is locally
  equivalent to $\exp(i\,\frac{\pi}{4}\,XX)$ via single-qubit $Z$
  rotations (same Weyl class $(\pi/4,0,0)$). Achieves depth 2
  from $\ket{00}$ with appropriate $A$.

\item \textbf{$ZZ$ axis}: $\exp(i\,\frac{\pi}{4}\,ZZ)$ is diagonal
  in the computational basis: it applies phases to $\ket{00}$,
  $\ket{01}$, $\ket{10}$, $\ket{11}$ but does not mix them.
  Consequently $\exp(i\,\frac{\pi}{4}\,ZZ)\cdot(A\otimes I)\ket{00}$
  has concurrence $\mathcal{C} = 0$ for any single-qubit $A$,
  since $A\otimes I$ maps $\ket{00}$ to a state in the
  $\{\ket{00}, \ket{10}\}$ subspace, and $ZZ$ cannot mix
  the two-qubit amplitudes. Producing $\mathcal{C}=1$
  requires prior superposition on \emph{both} qubits,
  raising the minimum depth to $d = 3$: $(H\otimes H) \cdot ZZ$.
  (Numerically verified: $ZZ\cdot(A\otimes I)\ket{00}$ achieves
  $\mathcal{C}=0$ for all single-qubit $A$; $(H\otimes H)\cdot ZZ$
  achieves $\mathcal{C} \approx 1.000$ at depth 3.)

\end{itemize}

\medskip
\noindent\textbf{Step 3: Rank $\geq 2$ gates.}
All rank-$\geq 2$ gates with $c_1 = \pi/4$ are perfect entanglers
but require at least two non-trivial single-qubit preparation gates
to achieve $\mathcal{C}=1$ from $\ket{00}$, giving depth $\geq 3$.
Rank-$\geq 2$ gates with $c_1 < \pi/4$ (e.g., $\sqrt{i\mathrm{SWAP}}$)
require two applications to achieve maximal entanglement, giving
depth $\geq 4$. Both cases are eliminated by T4.

\medskip
\noindent\textbf{Step 4: T1 compliance.}
The $XX$/$YY$ class (CNOT Weyl class) is directional:
$\mathrm{CNOT}(a, b) \neq \mathrm{CNOT}(b, a)$ (control and target
are not interchangeable). Numerically: $\mathrm{CNOT}\ket{01} = \ket{01}$
while $\mathrm{CNOT}_{\mathrm{rev}}\ket{01} = \ket{11}$. T1 satisfied.

\medskip
\noindent\textbf{Step 5: T3 compliance.}
Individual CNOT applications are unitary (reversible). T3, however,
applies to the \emph{accumulated recursion}: after the fold reaches
termination depth $D$, the fixed-point state $\ket{\Psi_0}$ encodes
an irreversible classical record (thermodynamic entropy production at
UV scale). The CNOT class generates this irreversibility at termination
depth consistently with T3; no additional constraint from T3 eliminates
the CNOT class.

\medskip
\noindent\textbf{Conclusion.}
Only the $XX$/$YY$ Weyl class at $(\pi/4, 0, 0)$ achieves
minimum depth $d^* = 2$ from $\ket{00}$ while satisfying T1--T3.
Since $XX$ and $YY$ at $c_1 = \pi/4$ are locally equivalent
(same Makhlin invariants, same Weyl class), there is exactly one
equivalence class satisfying all four constraints: the CNOT class.
\qed
\end{proof}

% ── Corollary ─────────────────────────────────────────────────
\begin{corollary}[Canonical Form]
\label{cor:H_CNOT}
The canonical minimum-depth realization of the CNOT class from
$\ket{00}$ is $F_{\mathrm{circuit}} = \mathrm{CNOT} \cdot (H \otimes I)$:
\begin{equation}
    F_{\mathrm{circuit}}\ket{00}
    = \mathrm{CNOT}\cdot\frac{\ket{0}+\ket{1}}{\sqrt{2}}\otimes\ket{0}
    = \frac{\ket{00}+\ket{11}}{\sqrt{2}}
    = \ket{\Phi^+}.
    \label{eq:canonical_form}
\end{equation}
The Hadamard $H$ is the unique single-qubit gate mapping $\ket{0}$
to the uniform superposition $\ket{+}$; combined with CNOT, it
achieves $\mathcal{C} = 1$ in exactly 2 gates.
\end{corollary}

% ── Physical interpretation ────────────────────────────────────
\subsubsection*{Physical Interpretation}

Theorem~\ref{thm:CNOT_uniqueness} sharpens the claim of
Section~\ref{sec:bell_derivation}. The Bell state $\ket{\Phi^+}$
and the CHSH violation $\langle\mathrm{CHSH}\rangle = 2\sqrt{2}$
are not artifacts of a convenient circuit choice. The four constraints
T1--T4 uniquely force the CNOT Weyl class; any other realization
either fails to generate entanglement (T2), is symmetric (T1), or
requires more gates than the fold's information-maximization primitive
would select (T4). The circuit $F_{\mathrm{circuit}}$ is not the
minimal realization \emph{we chose}: it is the minimal realization
\emph{that exists}, given the fold's structural requirements.

% ── Computational verification note ───────────────────────────
\subsubsection*{Computational Verification}

All numerical claims were verified using Python (NumPy, SciPy) on
2026-05-02. Verification script:
\texttt{UFUU\_Bell\_T4\_verification\_20260502.py} (this folder).
Key results:
\begin{itemize}
    \item $\mathcal{C}(F_{\mathrm{circuit}}\ket{00}) = 1.000000$.
    \item $ZZ\cdot(A\otimes I)\ket{00}$: $\mathcal{C} = 0$ for all $A \in SU(2)$.
    \item $(H\otimes H)\cdot ZZ\ket{00}$: $\mathcal{C} \approx 1.000$ (depth 3).
    \item T1: $\mathrm{CNOT}\ket{01} \neq \mathrm{CNOT}_{\mathrm{rev}}\ket{01}$.
    \item $XX$ and $YY$ at $\pi/4$: identical Makhlin invariants (same Weyl class).
\end{itemize}

% ── References ────────────────────────────────────────────────
% Add to main bibliography if not already present:
%
% \bibitem{Zhang2003}
%   Zhang J, Vala J, Sastry S, Whaley K B (2003).
%   Geometric theory of nonlocal two-qubit operations.
%   \textit{Physical Review A} \textbf{67}, 042313.
%
% \bibitem{Childs2003}
%   Childs A M, Haselgrove H L, Nielsen M A (2003).
%   Lower bounds on the complexity of simulating quantum gates.
%   \textit{Physical Review A} \textbf{68}, 052311.
%
% \bibitem{Makhlin2002}
%   Makhlin Y (2002).
%   Nonlocal properties of two-qubit gates and mixed states,
%   and the optimization of quantum computations.
%   \textit{Quantum Information Processing} \textbf{1}, 243--252.

% ── End standalone ─────────────────────────────────────────────
% \end{document}

%   or copy Section 3.3 verbatim into the main .tex file.
%
% Correction note (2026-05-02):
%   Earlier draft incorrectly claimed ZZ at pi/4 has zero
%   concurrence. ZZ IS a perfect entangler but requires circuit
%   depth 3 from |00> (two single-qubit prep gates + ZZ).
%   T4 is therefore correctly stated as minimum CIRCUIT DEPTH
%   from product input |00>, not minimum interaction rank alone.
%   The uniqueness result is unchanged; only the proof path differs.
% ============================================================

% ── If compiling standalone, uncomment the block below ──────
% \documentclass[12pt]{article}
% \usepackage{amsmath,amssymb,amsthm}
% \usepackage{physics}
% \usepackage{braket}
% \usepackage{hyperref}
% \newtheorem{theorem}{Theorem}
% \newtheorem{lemma}{Lemma}
% \newtheorem{definition}{Definition}
% \newtheorem{corollary}{Corollary}
% \begin{document}
% ── End standalone preamble ─────────────────────────────────

\subsection{Uniqueness of the Two-Qubit Realization (T4)}
\label{sec:T4_uniqueness}

Section~\ref{sec:fold_circuit} identified
$F_{\mathrm{circuit}} = \mathrm{CNOT} \cdot (H \otimes I)$
as the minimal unitary realization of the abstract fold operation $F$
in the two-qubit Hilbert space $\mathcal{H} = \mathbb{C}^2 \otimes \mathbb{C}^2$.
We now prove that this realization is \emph{unique} by introducing a
fourth structural constraint (T4) and showing that T1--T4 jointly select
the CNOT Weyl class as the only admissible realization, up to local unitaries.

% ── Weyl Chamber Background ──────────────────────────────────
\subsubsection*{Background: The Weyl Chamber}

Two-qubit gates in $SU(4)$ are classified up to local (single-qubit)
unitaries by their Weyl decomposition~\cite{Zhang2003}:
\begin{equation}
    U = (A_1 \otimes A_2)\,
        \exp\!\bigl[i(c_1\,XX + c_2\,YY + c_3\,ZZ)\bigr]\,
        (A_3 \otimes A_4),
    \label{eq:weyl_decomp}
\end{equation}
where $A_k \in SU(2)$, $XX = \sigma_x\otimes\sigma_x$,
$YY = \sigma_y\otimes\sigma_y$, $ZZ = \sigma_z\otimes\sigma_z$,
and $(c_1,c_2,c_3) \in [0,\pi/4]^3$ with $c_1 \geq c_2 \geq c_3 \geq 0$
are the Weyl coordinates. The interaction rank $\mathrm{rank}(U) =
\#\{k: c_k \neq 0\}$ counts the number of active interaction axes.
Key classes and their coordinates are given in Table~\ref{tab:weyl}.

\begin{table}[h]
\centering
\caption{Weyl coordinates and interaction rank for representative two-qubit gates.}
\label{tab:weyl}
\begin{tabular}{lcccc}
\hline
Gate class       & $c_1$    & $c_2$    & $c_3$    & Rank \\
\hline
Identity         & $0$      & $0$      & $0$      & 0 \\
CNOT             & $\pi/4$  & $0$      & $0$      & 1 \\
$\sqrt{i\mathrm{SWAP}}$ & $\pi/8$ & $\pi/8$ & $0$ & 2 \\
$i\mathrm{SWAP}$ & $\pi/4$  & $\pi/4$  & $0$      & 2 \\
B gate           & $\pi/4$  & $\pi/8$  & $0$      & 2 \\
$\sqrt{\mathrm{SWAP}}$ & $\pi/8$ & $\pi/8$ & $\pi/8$ & 3 \\
SWAP             & $\pi/4$  & $\pi/4$  & $\pi/4$  & 3 \\
\hline
\end{tabular}
\end{table}

% ── T4 Definition ────────────────────────────────────────────
\subsubsection*{The Fourth Constraint: Minimum Circuit Depth from Product Input (T4)}

\begin{definition}[T4 --- Minimum Circuit Depth]
\label{def:T4}
\textbf{(T4) Minimum circuit depth:}
The fold operation $F$ selects, from among all unitary realizations
satisfying T1--T3 and capable of producing maximal entanglement from
the unresolved binary product state $\ket{00}$, the realization requiring
the fewest elementary gates (single-qubit unitaries plus two-qubit
entangling gates) to achieve concurrence $\mathcal{C} = 1$ starting
from $\ket{00}$.
\end{definition}

The physical motivation is immediate. The recursive binary fold
$U = F(U,U)$ begins each step with two independent, unresolved binary
distinctions — represented as the product state $\ket{00}$. The
information-maximization primitive requires that maximal entanglement
be generated with minimum additional structure: the simplest circuit
that accomplishes this is the one the fold selects. T4 makes this
minimality requirement precise and measurable.

% ── Main Theorem ─────────────────────────────────────────────
\subsubsection*{Uniqueness Theorem}

\begin{theorem}[T1--T4 Uniquely Select the CNOT Class]
\label{thm:CNOT_uniqueness}
Among all two-qubit unitary gates, T1 (non-commutativity),
T2 (entanglement generation), T3 (partial irreversibility at
the recursive level), and T4 (minimum circuit depth from $\ket{00}$)
jointly select the CNOT Weyl class --- gates with Weyl coordinates
$(\pi/4, 0, 0)$ --- as the unique admissible realization of $F$,
up to local unitaries.
\end{theorem}

\begin{proof}
We establish the minimum achievable circuit depth (T4), then show
only the CNOT class achieves it while satisfying T1--T3.

\medskip
\noindent\textbf{Step 1: Lower bound on circuit depth.}
Any circuit producing maximal entanglement from $\ket{00}$ must
include at least one two-qubit entangling gate (since product
unitaries preserve the product structure of $\ket{00}$, violating T2)
and at least one single-qubit gate to create the superposition that
the entangling gate then correlates. The minimum circuit depth from
$\ket{00}$ to a maximally entangled state is therefore $d^* = 2$.

\medskip
\noindent\textbf{Step 2: Which rank-1 gates achieve $d^* = 2$ from $\ket{00}$?}
A depth-2 circuit from $\ket{00}$ has the form
$(A \otimes I) \cdot E$ or $(I \otimes A) \cdot E$ where $A \in SU(2)$
is a single-qubit gate and $E$ is a two-qubit entangling gate.
We require $\mathcal{C}(E\cdot(A\otimes I)\ket{00}) = 1$.

For the three single-axis interaction classes at $c = \pi/4$:

\begin{itemize}

\item \textbf{$XX$ axis}: $\exp(i\,\frac{\pi}{4}\,XX)$ with $A = H$
  (Hadamard) gives $\mathrm{CNOT}\cdot(H\otimes I)\ket{00} = \ket{\Phi^+}$,
  concurrence $= 1$. Depth 2 achieved. (Numerically verified:
  $\mathcal{C} = 1.000000$.)

\item \textbf{$YY$ axis}: $\exp(i\,\frac{\pi}{4}\,YY)$ is locally
  equivalent to $\exp(i\,\frac{\pi}{4}\,XX)$ via single-qubit $Z$
  rotations (same Weyl class $(\pi/4,0,0)$). Achieves depth 2
  from $\ket{00}$ with appropriate $A$.

\item \textbf{$ZZ$ axis}: $\exp(i\,\frac{\pi}{4}\,ZZ)$ is diagonal
  in the computational basis: it applies phases to $\ket{00}$,
  $\ket{01}$, $\ket{10}$, $\ket{11}$ but does not mix them.
  Consequently $\exp(i\,\frac{\pi}{4}\,ZZ)\cdot(A\otimes I)\ket{00}$
  has concurrence $\mathcal{C} = 0$ for any single-qubit $A$,
  since $A\otimes I$ maps $\ket{00}$ to a state in the
  $\{\ket{00}, \ket{10}\}$ subspace, and $ZZ$ cannot mix
  the two-qubit amplitudes. Producing $\mathcal{C}=1$
  requires prior superposition on \emph{both} qubits,
  raising the minimum depth to $d = 3$: $(H\otimes H) \cdot ZZ$.
  (Numerically verified: $ZZ\cdot(A\otimes I)\ket{00}$ achieves
  $\mathcal{C}=0$ for all single-qubit $A$; $(H\otimes H)\cdot ZZ$
  achieves $\mathcal{C} \approx 1.000$ at depth 3.)

\end{itemize}

\medskip
\noindent\textbf{Step 3: Rank $\geq 2$ gates.}
All rank-$\geq 2$ gates with $c_1 = \pi/4$ are perfect entanglers
but require at least two non-trivial single-qubit preparation gates
to achieve $\mathcal{C}=1$ from $\ket{00}$, giving depth $\geq 3$.
Rank-$\geq 2$ gates with $c_1 < \pi/4$ (e.g., $\sqrt{i\mathrm{SWAP}}$)
require two applications to achieve maximal entanglement, giving
depth $\geq 4$. Both cases are eliminated by T4.

\medskip
\noindent\textbf{Step 4: T1 compliance.}
The $XX$/$YY$ class (CNOT Weyl class) is directional:
$\mathrm{CNOT}(a, b) \neq \mathrm{CNOT}(b, a)$ (control and target
are not interchangeable). Numerically: $\mathrm{CNOT}\ket{01} = \ket{01}$
while $\mathrm{CNOT}_{\mathrm{rev}}\ket{01} = \ket{11}$. T1 satisfied.

\medskip
\noindent\textbf{Step 5: T3 compliance.}
Individual CNOT applications are unitary (reversible). T3, however,
applies to the \emph{accumulated recursion}: after the fold reaches
termination depth $D$, the fixed-point state $\ket{\Psi_0}$ encodes
an irreversible classical record (thermodynamic entropy production at
UV scale). The CNOT class generates this irreversibility at termination
depth consistently with T3; no additional constraint from T3 eliminates
the CNOT class.

\medskip
\noindent\textbf{Conclusion.}
Only the $XX$/$YY$ Weyl class at $(\pi/4, 0, 0)$ achieves
minimum depth $d^* = 2$ from $\ket{00}$ while satisfying T1--T3.
Since $XX$ and $YY$ at $c_1 = \pi/4$ are locally equivalent
(same Makhlin invariants, same Weyl class), there is exactly one
equivalence class satisfying all four constraints: the CNOT class.
\qed
\end{proof}

% ── Corollary ─────────────────────────────────────────────────
\begin{corollary}[Canonical Form]
\label{cor:H_CNOT}
The canonical minimum-depth realization of the CNOT class from
$\ket{00}$ is $F_{\mathrm{circuit}} = \mathrm{CNOT} \cdot (H \otimes I)$:
\begin{equation}
    F_{\mathrm{circuit}}\ket{00}
    = \mathrm{CNOT}\cdot\frac{\ket{0}+\ket{1}}{\sqrt{2}}\otimes\ket{0}
    = \frac{\ket{00}+\ket{11}}{\sqrt{2}}
    = \ket{\Phi^+}.
    \label{eq:canonical_form}
\end{equation}
The Hadamard $H$ is the unique single-qubit gate mapping $\ket{0}$
to the uniform superposition $\ket{+}$; combined with CNOT, it
achieves $\mathcal{C} = 1$ in exactly 2 gates.
\end{corollary}

% ── Physical interpretation ────────────────────────────────────
\subsubsection*{Physical Interpretation}

Theorem~\ref{thm:CNOT_uniqueness} sharpens the claim of
Section~\ref{sec:bell_derivation}. The Bell state $\ket{\Phi^+}$
and the CHSH violation $\langle\mathrm{CHSH}\rangle = 2\sqrt{2}$
are not artifacts of a convenient circuit choice. The four constraints
T1--T4 uniquely force the CNOT Weyl class; any other realization
either fails to generate entanglement (T2), is symmetric (T1), or
requires more gates than the fold's information-maximization primitive
would select (T4). The circuit $F_{\mathrm{circuit}}$ is not the
minimal realization \emph{we chose}: it is the minimal realization
\emph{that exists}, given the fold's structural requirements.

% ── Computational verification note ───────────────────────────
\subsubsection*{Computational Verification}

All numerical claims were verified using Python (NumPy, SciPy) on
2026-05-02. Verification script:
\texttt{UFUU\_Bell\_T4\_verification\_20260502.py} (this folder).
Key results:
\begin{itemize}
    \item $\mathcal{C}(F_{\mathrm{circuit}}\ket{00}) = 1.000000$.
    \item $ZZ\cdot(A\otimes I)\ket{00}$: $\mathcal{C} = 0$ for all $A \in SU(2)$.
    \item $(H\otimes H)\cdot ZZ\ket{00}$: $\mathcal{C} \approx 1.000$ (depth 3).
    \item T1: $\mathrm{CNOT}\ket{01} \neq \mathrm{CNOT}_{\mathrm{rev}}\ket{01}$.
    \item $XX$ and $YY$ at $\pi/4$: identical Makhlin invariants (same Weyl class).
\end{itemize}

% ── References ────────────────────────────────────────────────
% Add to main bibliography if not already present:
%
% \bibitem{Zhang2003}
%   Zhang J, Vala J, Sastry S, Whaley K B (2003).
%   Geometric theory of nonlocal two-qubit operations.
%   \textit{Physical Review A} \textbf{67}, 042313.
%
% \bibitem{Childs2003}
%   Childs A M, Haselgrove H L, Nielsen M A (2003).
%   Lower bounds on the complexity of simulating quantum gates.
%   \textit{Physical Review A} \textbf{68}, 052311.
%
% \bibitem{Makhlin2002}
%   Makhlin Y (2002).
%   Nonlocal properties of two-qubit gates and mixed states,
%   and the optimization of quantum computations.
%   \textit{Quantum Information Processing} \textbf{1}, 243--252.

% ── End standalone ─────────────────────────────────────────────
% \end{document}

%   or copy Section 3.3 verbatim into the main .tex file.
%
% Correction note (2026-05-02):
%   Earlier draft incorrectly claimed ZZ at pi/4 has zero
%   concurrence. ZZ IS a perfect entangler but requires circuit
%   depth 3 from |00> (two single-qubit prep gates + ZZ).
%   T4 is therefore correctly stated as minimum CIRCUIT DEPTH
%   from product input |00>, not minimum interaction rank alone.
%   The uniqueness result is unchanged; only the proof path differs.
% ============================================================

% ── If compiling standalone, uncomment the block below ──────
% \documentclass[12pt]{article}
% \usepackage{amsmath,amssymb,amsthm}
% \usepackage{physics}
% \usepackage{braket}
% \usepackage{hyperref}
% \newtheorem{theorem}{Theorem}
% \newtheorem{lemma}{Lemma}
% \newtheorem{definition}{Definition}
% \newtheorem{corollary}{Corollary}
% \begin{document}
% ── End standalone preamble ─────────────────────────────────

\subsection{Uniqueness of the Two-Qubit Realization (T4)}
\label{sec:T4_uniqueness}

Section~\ref{sec:fold_circuit} identified
$F_{\mathrm{circuit}} = \mathrm{CNOT} \cdot (H \otimes I)$
as the minimal unitary realization of the abstract fold operation $F$
in the two-qubit Hilbert space $\mathcal{H} = \mathbb{C}^2 \otimes \mathbb{C}^2$.
We now prove that this realization is \emph{unique} by introducing a
fourth structural constraint (T4) and showing that T1--T4 jointly select
the CNOT Weyl class as the only admissible realization, up to local unitaries.

% ── Weyl Chamber Background ──────────────────────────────────
\subsubsection*{Background: The Weyl Chamber}

Two-qubit gates in $SU(4)$ are classified up to local (single-qubit)
unitaries by their Weyl decomposition~\cite{Zhang2003}:
\begin{equation}
    U = (A_1 \otimes A_2)\,
        \exp\!\bigl[i(c_1\,XX + c_2\,YY + c_3\,ZZ)\bigr]\,
        (A_3 \otimes A_4),
    \label{eq:weyl_decomp}
\end{equation}
where $A_k \in SU(2)$, $XX = \sigma_x\otimes\sigma_x$,
$YY = \sigma_y\otimes\sigma_y$, $ZZ = \sigma_z\otimes\sigma_z$,
and $(c_1,c_2,c_3) \in [0,\pi/4]^3$ with $c_1 \geq c_2 \geq c_3 \geq 0$
are the Weyl coordinates. The interaction rank $\mathrm{rank}(U) =
\#\{k: c_k \neq 0\}$ counts the number of active interaction axes.
Key classes and their coordinates are given in Table~\ref{tab:weyl}.

\begin{table}[h]
\centering
\caption{Weyl coordinates and interaction rank for representative two-qubit gates.}
\label{tab:weyl}
\begin{tabular}{lcccc}
\hline
Gate class       & $c_1$    & $c_2$    & $c_3$    & Rank \\
\hline
Identity         & $0$      & $0$      & $0$      & 0 \\
CNOT             & $\pi/4$  & $0$      & $0$      & 1 \\
$\sqrt{i\mathrm{SWAP}}$ & $\pi/8$ & $\pi/8$ & $0$ & 2 \\
$i\mathrm{SWAP}$ & $\pi/4$  & $\pi/4$  & $0$      & 2 \\
B gate           & $\pi/4$  & $\pi/8$  & $0$      & 2 \\
$\sqrt{\mathrm{SWAP}}$ & $\pi/8$ & $\pi/8$ & $\pi/8$ & 3 \\
SWAP             & $\pi/4$  & $\pi/4$  & $\pi/4$  & 3 \\
\hline
\end{tabular}
\end{table}

% ── T4 Definition ────────────────────────────────────────────
\subsubsection*{The Fourth Constraint: Minimum Circuit Depth from Product Input (T4)}

\begin{definition}[T4 --- Minimum Circuit Depth]
\label{def:T4}
\textbf{(T4) Minimum circuit depth:}
The fold operation $F$ selects, from among all unitary realizations
satisfying T1--T3 and capable of producing maximal entanglement from
the unresolved binary product state $\ket{00}$, the realization requiring
the fewest elementary gates (single-qubit unitaries plus two-qubit
entangling gates) to achieve concurrence $\mathcal{C} = 1$ starting
from $\ket{00}$.
\end{definition}

The physical motivation is immediate. The recursive binary fold
$U = F(U,U)$ begins each step with two independent, unresolved binary
distinctions — represented as the product state $\ket{00}$. The
information-maximization primitive requires that maximal entanglement
be generated with minimum additional structure: the simplest circuit
that accomplishes this is the one the fold selects. T4 makes this
minimality requirement precise and measurable.

% ── Main Theorem ─────────────────────────────────────────────
\subsubsection*{Uniqueness Theorem}

\begin{theorem}[T1--T4 Uniquely Select the CNOT Class]
\label{thm:CNOT_uniqueness}
Among all two-qubit unitary gates, T1 (non-commutativity),
T2 (entanglement generation), T3 (partial irreversibility at
the recursive level), and T4 (minimum circuit depth from $\ket{00}$)
jointly select the CNOT Weyl class --- gates with Weyl coordinates
$(\pi/4, 0, 0)$ --- as the unique admissible realization of $F$,
up to local unitaries.
\end{theorem}

\begin{proof}
We establish the minimum achievable circuit depth (T4), then show
only the CNOT class achieves it while satisfying T1--T3.

\medskip
\noindent\textbf{Step 1: Lower bound on circuit depth.}
Any circuit producing maximal entanglement from $\ket{00}$ must
include at least one two-qubit entangling gate (since product
unitaries preserve the product structure of $\ket{00}$, violating T2)
and at least one single-qubit gate to create the superposition that
the entangling gate then correlates. The minimum circuit depth from
$\ket{00}$ to a maximally entangled state is therefore $d^* = 2$.

\medskip
\noindent\textbf{Step 2: Which rank-1 gates achieve $d^* = 2$ from $\ket{00}$?}
A depth-2 circuit from $\ket{00}$ has the form
$(A \otimes I) \cdot E$ or $(I \otimes A) \cdot E$ where $A \in SU(2)$
is a single-qubit gate and $E$ is a two-qubit entangling gate.
We require $\mathcal{C}(E\cdot(A\otimes I)\ket{00}) = 1$.

For the three single-axis interaction classes at $c = \pi/4$:

\begin{itemize}

\item \textbf{$XX$ axis}: $\exp(i\,\frac{\pi}{4}\,XX)$ with $A = H$
  (Hadamard) gives $\mathrm{CNOT}\cdot(H\otimes I)\ket{00} = \ket{\Phi^+}$,
  concurrence $= 1$. Depth 2 achieved. (Numerically verified:
  $\mathcal{C} = 1.000000$.)

\item \textbf{$YY$ axis}: $\exp(i\,\frac{\pi}{4}\,YY)$ is locally
  equivalent to $\exp(i\,\frac{\pi}{4}\,XX)$ via single-qubit $Z$
  rotations (same Weyl class $(\pi/4,0,0)$). Achieves depth 2
  from $\ket{00}$ with appropriate $A$.

\item \textbf{$ZZ$ axis}: $\exp(i\,\frac{\pi}{4}\,ZZ)$ is diagonal
  in the computational basis: it applies phases to $\ket{00}$,
  $\ket{01}$, $\ket{10}$, $\ket{11}$ but does not mix them.
  Consequently $\exp(i\,\frac{\pi}{4}\,ZZ)\cdot(A\otimes I)\ket{00}$
  has concurrence $\mathcal{C} = 0$ for any single-qubit $A$,
  since $A\otimes I$ maps $\ket{00}$ to a state in the
  $\{\ket{00}, \ket{10}\}$ subspace, and $ZZ$ cannot mix
  the two-qubit amplitudes. Producing $\mathcal{C}=1$
  requires prior superposition on \emph{both} qubits,
  raising the minimum depth to $d = 3$: $(H\otimes H) \cdot ZZ$.
  (Numerically verified: $ZZ\cdot(A\otimes I)\ket{00}$ achieves
  $\mathcal{C}=0$ for all single-qubit $A$; $(H\otimes H)\cdot ZZ$
  achieves $\mathcal{C} \approx 1.000$ at depth 3.)

\end{itemize}

\medskip
\noindent\textbf{Step 3: Rank $\geq 2$ gates.}
All rank-$\geq 2$ gates with $c_1 = \pi/4$ are perfect entanglers
but require at least two non-trivial single-qubit preparation gates
to achieve $\mathcal{C}=1$ from $\ket{00}$, giving depth $\geq 3$.
Rank-$\geq 2$ gates with $c_1 < \pi/4$ (e.g., $\sqrt{i\mathrm{SWAP}}$)
require two applications to achieve maximal entanglement, giving
depth $\geq 4$. Both cases are eliminated by T4.

\medskip
\noindent\textbf{Step 4: T1 compliance.}
The $XX$/$YY$ class (CNOT Weyl class) is directional:
$\mathrm{CNOT}(a, b) \neq \mathrm{CNOT}(b, a)$ (control and target
are not interchangeable). Numerically: $\mathrm{CNOT}\ket{01} = \ket{01}$
while $\mathrm{CNOT}_{\mathrm{rev}}\ket{01} = \ket{11}$. T1 satisfied.

\medskip
\noindent\textbf{Step 5: T3 compliance.}
Individual CNOT applications are unitary (reversible). T3, however,
applies to the \emph{accumulated recursion}: after the fold reaches
termination depth $D$, the fixed-point state $\ket{\Psi_0}$ encodes
an irreversible classical record (thermodynamic entropy production at
UV scale). The CNOT class generates this irreversibility at termination
depth consistently with T3; no additional constraint from T3 eliminates
the CNOT class.

\medskip
\noindent\textbf{Conclusion.}
Only the $XX$/$YY$ Weyl class at $(\pi/4, 0, 0)$ achieves
minimum depth $d^* = 2$ from $\ket{00}$ while satisfying T1--T3.
Since $XX$ and $YY$ at $c_1 = \pi/4$ are locally equivalent
(same Makhlin invariants, same Weyl class), there is exactly one
equivalence class satisfying all four constraints: the CNOT class.
\qed
\end{proof}

% ── Corollary ─────────────────────────────────────────────────
\begin{corollary}[Canonical Form]
\label{cor:H_CNOT}
The canonical minimum-depth realization of the CNOT class from
$\ket{00}$ is $F_{\mathrm{circuit}} = \mathrm{CNOT} \cdot (H \otimes I)$:
\begin{equation}
    F_{\mathrm{circuit}}\ket{00}
    = \mathrm{CNOT}\cdot\frac{\ket{0}+\ket{1}}{\sqrt{2}}\otimes\ket{0}
    = \frac{\ket{00}+\ket{11}}{\sqrt{2}}
    = \ket{\Phi^+}.
    \label{eq:canonical_form}
\end{equation}
The Hadamard $H$ is the unique single-qubit gate mapping $\ket{0}$
to the uniform superposition $\ket{+}$; combined with CNOT, it
achieves $\mathcal{C} = 1$ in exactly 2 gates.
\end{corollary}

% ── Physical interpretation ────────────────────────────────────
\subsubsection*{Physical Interpretation}

Theorem~\ref{thm:CNOT_uniqueness} sharpens the claim of
Section~\ref{sec:bell_derivation}. The Bell state $\ket{\Phi^+}$
and the CHSH violation $\langle\mathrm{CHSH}\rangle = 2\sqrt{2}$
are not artifacts of a convenient circuit choice. The four constraints
T1--T4 uniquely force the CNOT Weyl class; any other realization
either fails to generate entanglement (T2), is symmetric (T1), or
requires more gates than the fold's information-maximization primitive
would select (T4). The circuit $F_{\mathrm{circuit}}$ is not the
minimal realization \emph{we chose}: it is the minimal realization
\emph{that exists}, given the fold's structural requirements.

% ── Computational verification note ───────────────────────────
\subsubsection*{Computational Verification}

All numerical claims were verified using Python (NumPy, SciPy) on
2026-05-02. Verification script:
\texttt{UFUU\_Bell\_T4\_verification\_20260502.py} (this folder).
Key results:
\begin{itemize}
    \item $\mathcal{C}(F_{\mathrm{circuit}}\ket{00}) = 1.000000$.
    \item $ZZ\cdot(A\otimes I)\ket{00}$: $\mathcal{C} = 0$ for all $A \in SU(2)$.
    \item $(H\otimes H)\cdot ZZ\ket{00}$: $\mathcal{C} \approx 1.000$ (depth 3).
    \item T1: $\mathrm{CNOT}\ket{01} \neq \mathrm{CNOT}_{\mathrm{rev}}\ket{01}$.
    \item $XX$ and $YY$ at $\pi/4$: identical Makhlin invariants (same Weyl class).
\end{itemize}

% ── References ────────────────────────────────────────────────
% Add to main bibliography if not already present:
%
% \bibitem{Zhang2003}
%   Zhang J, Vala J, Sastry S, Whaley K B (2003).
%   Geometric theory of nonlocal two-qubit operations.
%   \textit{Physical Review A} \textbf{67}, 042313.
%
% \bibitem{Childs2003}
%   Childs A M, Haselgrove H L, Nielsen M A (2003).
%   Lower bounds on the complexity of simulating quantum gates.
%   \textit{Physical Review A} \textbf{68}, 052311.
%
% \bibitem{Makhlin2002}
%   Makhlin Y (2002).
%   Nonlocal properties of two-qubit gates and mixed states,
%   and the optimization of quantum computations.
%   \textit{Quantum Information Processing} \textbf{1}, 243--252.

% ── End standalone ─────────────────────────────────────────────
% \end{document}

%   or copy Section 3.3 verbatim into the main .tex file.
%
% Correction note (2026-05-02):
%   Earlier draft incorrectly claimed ZZ at pi/4 has zero
%   concurrence. ZZ IS a perfect entangler but requires circuit
%   depth 3 from |00> (two single-qubit prep gates + ZZ).
%   T4 is therefore correctly stated as minimum CIRCUIT DEPTH
%   from product input |00>, not minimum interaction rank alone.
%   The uniqueness result is unchanged; only the proof path differs.
% ============================================================

% ── If compiling standalone, uncomment the block below ──────
% \documentclass[12pt]{article}
% \usepackage{amsmath,amssymb,amsthm}
% \usepackage{physics}
% \usepackage{braket}
% \usepackage{hyperref}
% \newtheorem{theorem}{Theorem}
% \newtheorem{lemma}{Lemma}
% \newtheorem{definition}{Definition}
% \newtheorem{corollary}{Corollary}
% \begin{document}
% ── End standalone preamble ─────────────────────────────────

\subsection{Uniqueness of the Two-Qubit Realization (T4)}
\label{sec:T4_uniqueness}

Section~\ref{sec:fold_circuit} identified
$F_{\mathrm{circuit}} = \mathrm{CNOT} \cdot (H \otimes I)$
as the minimal unitary realization of the abstract fold operation $F$
in the two-qubit Hilbert space $\mathcal{H} = \mathbb{C}^2 \otimes \mathbb{C}^2$.
We now prove that this realization is \emph{unique} by introducing a
fourth structural constraint (T4) and showing that T1--T4 jointly select
the CNOT Weyl class as the only admissible realization, up to local unitaries.

% ── Weyl Chamber Background ──────────────────────────────────
\subsubsection*{Background: The Weyl Chamber}

Two-qubit gates in $SU(4)$ are classified up to local (single-qubit)
unitaries by their Weyl decomposition~\cite{Zhang2003}:
\begin{equation}
    U = (A_1 \otimes A_2)\,
        \exp\!\bigl[i(c_1\,XX + c_2\,YY + c_3\,ZZ)\bigr]\,
        (A_3 \otimes A_4),
    \label{eq:weyl_decomp}
\end{equation}
where $A_k \in SU(2)$, $XX = \sigma_x\otimes\sigma_x$,
$YY = \sigma_y\otimes\sigma_y$, $ZZ = \sigma_z\otimes\sigma_z$,
and $(c_1,c_2,c_3) \in [0,\pi/4]^3$ with $c_1 \geq c_2 \geq c_3 \geq 0$
are the Weyl coordinates. The interaction rank $\mathrm{rank}(U) =
\#\{k: c_k \neq 0\}$ counts the number of active interaction axes.
Key classes and their coordinates are given in Table~\ref{tab:weyl}.

\begin{table}[h]
\centering
\caption{Weyl coordinates and interaction rank for representative two-qubit gates.}
\label{tab:weyl}
\begin{tabular}{lcccc}
\hline
Gate class       & $c_1$    & $c_2$    & $c_3$    & Rank \\
\hline
Identity         & $0$      & $0$      & $0$      & 0 \\
CNOT             & $\pi/4$  & $0$      & $0$      & 1 \\
$\sqrt{i\mathrm{SWAP}}$ & $\pi/8$ & $\pi/8$ & $0$ & 2 \\
$i\mathrm{SWAP}$ & $\pi/4$  & $\pi/4$  & $0$      & 2 \\
B gate           & $\pi/4$  & $\pi/8$  & $0$      & 2 \\
$\sqrt{\mathrm{SWAP}}$ & $\pi/8$ & $\pi/8$ & $\pi/8$ & 3 \\
SWAP             & $\pi/4$  & $\pi/4$  & $\pi/4$  & 3 \\
\hline
\end{tabular}
\end{table}

% ── T4 Definition ────────────────────────────────────────────
\subsubsection*{The Fourth Constraint: Minimum Circuit Depth from Product Input (T4)}

\begin{definition}[T4 --- Minimum Circuit Depth]
\label{def:T4}
\textbf{(T4) Minimum circuit depth:}
The fold operation $F$ selects, from among all unitary realizations
satisfying T1--T3 and capable of producing maximal entanglement from
the unresolved binary product state $\ket{00}$, the realization requiring
the fewest elementary gates (single-qubit unitaries plus two-qubit
entangling gates) to achieve concurrence $\mathcal{C} = 1$ starting
from $\ket{00}$.
\end{definition}

The physical motivation is immediate. The recursive binary fold
$U = F(U,U)$ begins each step with two independent, unresolved binary
distinctions — represented as the product state $\ket{00}$. The
information-maximization primitive requires that maximal entanglement
be generated with minimum additional structure: the simplest circuit
that accomplishes this is the one the fold selects. T4 makes this
minimality requirement precise and measurable.

% ── Main Theorem ─────────────────────────────────────────────
\subsubsection*{Uniqueness Theorem}

\begin{theorem}[T1--T4 Uniquely Select the CNOT Class]
\label{thm:CNOT_uniqueness}
Among all two-qubit unitary gates, T1 (non-commutativity),
T2 (entanglement generation), T3 (partial irreversibility at
the recursive level), and T4 (minimum circuit depth from $\ket{00}$)
jointly select the CNOT Weyl class --- gates with Weyl coordinates
$(\pi/4, 0, 0)$ --- as the unique admissible realization of $F$,
up to local unitaries.
\end{theorem}

\begin{proof}
We establish the minimum achievable circuit depth (T4), then show
only the CNOT class achieves it while satisfying T1--T3.

\medskip
\noindent\textbf{Step 1: Lower bound on circuit depth.}
Any circuit producing maximal entanglement from $\ket{00}$ must
include at least one two-qubit entangling gate (since product
unitaries preserve the product structure of $\ket{00}$, violating T2)
and at least one single-qubit gate to create the superposition that
the entangling gate then correlates. The minimum circuit depth from
$\ket{00}$ to a maximally entangled state is therefore $d^* = 2$.

\medskip
\noindent\textbf{Step 2: Which rank-1 gates achieve $d^* = 2$ from $\ket{00}$?}
A depth-2 circuit from $\ket{00}$ has the form
$(A \otimes I) \cdot E$ or $(I \otimes A) \cdot E$ where $A \in SU(2)$
is a single-qubit gate and $E$ is a two-qubit entangling gate.
We require $\mathcal{C}(E\cdot(A\otimes I)\ket{00}) = 1$.

For the three single-axis interaction classes at $c = \pi/4$:

\begin{itemize}

\item \textbf{$XX$ axis}: $\exp(i\,\frac{\pi}{4}\,XX)$ with $A = H$
  (Hadamard) gives $\mathrm{CNOT}\cdot(H\otimes I)\ket{00} = \ket{\Phi^+}$,
  concurrence $= 1$. Depth 2 achieved. (Numerically verified:
  $\mathcal{C} = 1.000000$.)

\item \textbf{$YY$ axis}: $\exp(i\,\frac{\pi}{4}\,YY)$ is locally
  equivalent to $\exp(i\,\frac{\pi}{4}\,XX)$ via single-qubit $Z$
  rotations (same Weyl class $(\pi/4,0,0)$). Achieves depth 2
  from $\ket{00}$ with appropriate $A$.

\item \textbf{$ZZ$ axis}: $\exp(i\,\frac{\pi}{4}\,ZZ)$ is diagonal
  in the computational basis: it applies phases to $\ket{00}$,
  $\ket{01}$, $\ket{10}$, $\ket{11}$ but does not mix them.
  Consequently $\exp(i\,\frac{\pi}{4}\,ZZ)\cdot(A\otimes I)\ket{00}$
  has concurrence $\mathcal{C} = 0$ for any single-qubit $A$,
  since $A\otimes I$ maps $\ket{00}$ to a state in the
  $\{\ket{00}, \ket{10}\}$ subspace, and $ZZ$ cannot mix
  the two-qubit amplitudes. Producing $\mathcal{C}=1$
  requires prior superposition on \emph{both} qubits,
  raising the minimum depth to $d = 3$: $(H\otimes H) \cdot ZZ$.
  (Numerically verified: $ZZ\cdot(A\otimes I)\ket{00}$ achieves
  $\mathcal{C}=0$ for all single-qubit $A$; $(H\otimes H)\cdot ZZ$
  achieves $\mathcal{C} \approx 1.000$ at depth 3.)

\end{itemize}

\medskip
\noindent\textbf{Step 3: Rank $\geq 2$ gates.}
All rank-$\geq 2$ gates with $c_1 = \pi/4$ are perfect entanglers
but require at least two non-trivial single-qubit preparation gates
to achieve $\mathcal{C}=1$ from $\ket{00}$, giving depth $\geq 3$.
Rank-$\geq 2$ gates with $c_1 < \pi/4$ (e.g., $\sqrt{i\mathrm{SWAP}}$)
require two applications to achieve maximal entanglement, giving
depth $\geq 4$. Both cases are eliminated by T4.

\medskip
\noindent\textbf{Step 4: T1 compliance.}
The $XX$/$YY$ class (CNOT Weyl class) is directional:
$\mathrm{CNOT}(a, b) \neq \mathrm{CNOT}(b, a)$ (control and target
are not interchangeable). Numerically: $\mathrm{CNOT}\ket{01} = \ket{01}$
while $\mathrm{CNOT}_{\mathrm{rev}}\ket{01} = \ket{11}$. T1 satisfied.

\medskip
\noindent\textbf{Step 5: T3 compliance.}
Individual CNOT applications are unitary (reversible). T3, however,
applies to the \emph{accumulated recursion}: after the fold reaches
termination depth $D$, the fixed-point state $\ket{\Psi_0}$ encodes
an irreversible classical record (thermodynamic entropy production at
UV scale). The CNOT class generates this irreversibility at termination
depth consistently with T3; no additional constraint from T3 eliminates
the CNOT class.

\medskip
\noindent\textbf{Conclusion.}
Only the $XX$/$YY$ Weyl class at $(\pi/4, 0, 0)$ achieves
minimum depth $d^* = 2$ from $\ket{00}$ while satisfying T1--T3.
Since $XX$ and $YY$ at $c_1 = \pi/4$ are locally equivalent
(same Makhlin invariants, same Weyl class), there is exactly one
equivalence class satisfying all four constraints: the CNOT class.
\qed
\end{proof}

% ── Corollary ─────────────────────────────────────────────────
\begin{corollary}[Canonical Form]
\label{cor:H_CNOT}
The canonical minimum-depth realization of the CNOT class from
$\ket{00}$ is $F_{\mathrm{circuit}} = \mathrm{CNOT} \cdot (H \otimes I)$:
\begin{equation}
    F_{\mathrm{circuit}}\ket{00}
    = \mathrm{CNOT}\cdot\frac{\ket{0}+\ket{1}}{\sqrt{2}}\otimes\ket{0}
    = \frac{\ket{00}+\ket{11}}{\sqrt{2}}
    = \ket{\Phi^+}.
    \label{eq:canonical_form}
\end{equation}
The Hadamard $H$ is the unique single-qubit gate mapping $\ket{0}$
to the uniform superposition $\ket{+}$; combined with CNOT, it
achieves $\mathcal{C} = 1$ in exactly 2 gates.
\end{corollary}

% ── Physical interpretation ────────────────────────────────────
\subsubsection*{Physical Interpretation}

Theorem~\ref{thm:CNOT_uniqueness} sharpens the claim of
Section~\ref{sec:bell_derivation}. The Bell state $\ket{\Phi^+}$
and the CHSH violation $\langle\mathrm{CHSH}\rangle = 2\sqrt{2}$
are not artifacts of a convenient circuit choice. The four constraints
T1--T4 uniquely force the CNOT Weyl class; any other realization
either fails to generate entanglement (T2), is symmetric (T1), or
requires more gates than the fold's information-maximization primitive
would select (T4). The circuit $F_{\mathrm{circuit}}$ is not the
minimal realization \emph{we chose}: it is the minimal realization
\emph{that exists}, given the fold's structural requirements.

% ── Computational verification note ───────────────────────────
\subsubsection*{Computational Verification}

All numerical claims were verified using Python (NumPy, SciPy) on
2026-05-02. Verification script:
\texttt{UFUU\_Bell\_T4\_verification\_20260502.py} (this folder).
Key results:
\begin{itemize}
    \item $\mathcal{C}(F_{\mathrm{circuit}}\ket{00}) = 1.000000$.
    \item $ZZ\cdot(A\otimes I)\ket{00}$: $\mathcal{C} = 0$ for all $A \in SU(2)$.
    \item $(H\otimes H)\cdot ZZ\ket{00}$: $\mathcal{C} \approx 1.000$ (depth 3).
    \item T1: $\mathrm{CNOT}\ket{01} \neq \mathrm{CNOT}_{\mathrm{rev}}\ket{01}$.
    \item $XX$ and $YY$ at $\pi/4$: identical Makhlin invariants (same Weyl class).
\end{itemize}

% ── References ────────────────────────────────────────────────
% Add to main bibliography if not already present:
%
% \bibitem{Zhang2003}
%   Zhang J, Vala J, Sastry S, Whaley K B (2003).
%   Geometric theory of nonlocal two-qubit operations.
%   \textit{Physical Review A} \textbf{67}, 042313.
%
% \bibitem{Childs2003}
%   Childs A M, Haselgrove H L, Nielsen M A (2003).
%   Lower bounds on the complexity of simulating quantum gates.
%   \textit{Physical Review A} \textbf{68}, 052311.
%
% \bibitem{Makhlin2002}
%   Makhlin Y (2002).
%   Nonlocal properties of two-qubit gates and mixed states,
%   and the optimization of quantum computations.
%   \textit{Quantum Information Processing} \textbf{1}, 243--252.

% ── End standalone ─────────────────────────────────────────────
% \end{document}
